% =====================================================================
%  GMN3 §7.2–7.5  扭转理论在黎曼面上的紧化 / 模空间 / 来自六维的线算子 / 线算子的编目
%  原文：Lamination-Paper.tex 4379--4685 行；PDF p.67--72（印刷 p.66--71）
%  编号公式 (7.6)--(7.14) 共 9 条（§7 共 14 条，§7.1 已用去 (7.1)--(7.5)）；
%  块内引用用 \eqref；跨块引用（式 (7.4)、(2.2) 及各节号）暂写原文编号并加 % TODO
% =====================================================================

\subsection{扭转理论在黎曼面上的紧化}\label{subsec:Twisted-C}

我们要考虑把 $(2,0)$ 理论缠绕在黎曼面 $C$ 上，这正是 \cite{Gaiotto:2009hg} 所用的构造。如 \cite{Gaiotto:2009hg} 第 3.1.2 节所述，为使这一紧化保持 $d=4,\CN=2$ 超对称，我们必须对理论作扭转。紧化 $\IR^{1,3}\times C$ 把局部洛伦兹对称性破缺为 $so(1,3)\oplus so(2)_C$。随后我们显式地破缺 $R$ 对称性 $so(2)'\oplus so(3)'\subset so(5)$，并对理论作扭转，使 $C$ 上的自旋联络与对角子代数 $so(2)_d\subset so(2)_C\oplus so(2)'$ 耦合。有八个超荷存活下来，它们在四维中构成一个 $\CN=2$ 超对称代数。$so(3)'$ 等同于 $d=4,\CN=2$ 的 $R$ 对称性，$so(2)_d$ 则是 $d=4,\CN=2$ 的反常 $u(1)_R$ 对称性。指标为 $d_k$ 的那族手征主算子中，含有一个 $so(2)_d$ 荷最大的算子，记作 $\CO_k$，于是 vev $\langle\CO_k\rangle$ 成为 $C$ 上的全纯 $k$-微分，而四维理论的库仑分支由这些算子的 vev 参数化（见 \cite{Gaiotto:2009hg} 式 (3.5)）：
\begin{equation}
\CM(\IR^{1,3}\times C) = \CB = \oplus H^0(C, K^{\otimes d_k})
\end{equation}
与 \cite{Gaiotto:2009hg} 中一样，重要的是引入安置在 $C$ 的穿刺点处的余维 2 缺陷。这些缺陷不再破缺任何超对称。它们的作用是在 $\langle\CO_k\rangle$ 中制造出极点。四维红外理论由 $C$ 的选取以及 $\langle\CO_k\rangle$ 在 $C$ 上各缺陷点处的奇异性决定。规范耦合的模空间与带穿刺点的黎曼面 $C$ 的复结构模空间重合。

\subsection{模空间}

如 \cite{Gaiotto:2009hg} 中详细解释的那样，真空模空间
\begin{equation}
\CM = \CM(\IR^{1,2}\times S^1\times C)
\end{equation}
可以等同于 $C$ 上 Hitchin 方程之解的模空间，这些解在缺陷处具有确定的奇异性。这些方程是关于联络 $A$ 与 1-形式 $\varphi$ 的方程，其内容是：复联络
\begin{equation} \label{eq:flat-conn}
\CA(\zeta) = R\frac{\varphi}{\zeta} + A + R\zeta\bar\varphi
\end{equation}
对一切 $\zeta\in\IC^\times$ 都平坦。事实上，在其复结构 $J^{(\zeta)}$ 下，模空间 $\CM$ 恰好等同于 $C$ 上平坦 $G_\IC$-联络 $\CA(\zeta)$ 的模空间，而这些联络在 $C$ 的穿刺点处具有奇异性。规范群 $G$ 的李代数为 $\lieg$。

穿刺点处奇异性的性质，取决于我们在 $C$ 上插入的具体缺陷。最简单的可能是「正则穿刺点」，它在 $A$ 与 $\varphi$ 中产生单极点，从而在 $\CA(\zeta)$ 中产生正则奇点。这样的穿刺点确定 $G_\IC$ 中一个（依赖于 $\zeta$ 的）共轭类；$\CA(\zeta)$ 的单值性被限制在该共轭类中取值。人们也可以考虑「非正则穿刺点」，它们产生更为复杂的非正则奇点。\zy{「正则/非正则穿刺点」是按其在 $\CA(\zeta)$ 中所产生的正则/非正则奇点命名的，与穿刺点本身无关。}

我们可以更具体地描述 $\CM$。若只考虑正则穿刺点 $\CP_i$，则平坦联络在模去规范变换之后完全由其单值表示捕获：于是 $\CM$ 是 $\pi_1(C\setminus\{\CP_i\})$ 到 $G_\IC$ 的表示所构成的空间，其中绕穿刺点的单值性的共轭类固定。\footnote{我们说「一个表示空间」而不说「\emph{那个}表示空间」，是因为涉及离散商的微妙之处，这一点稍后会讨论。} 再具体一些：取基点 $x\in C$，取绕 $C$ 上 $A$ 圈与 $B$ 圈的基于基点的环路，以及绕穿刺点 $\CP_i$ 的基于基点的环路；于是我们研究的是满足
\begin{equation}
M_1 \cdots M_n = \prod_{i=1}^{g} A_i B_i A_i^{-1} B_i^{-1},
\end{equation}
的多元组 $\{A_1,\dots,A_g,B_1,\dots,B_g,M_1,\dots,M_n\}\in G_\IC^{2g+n}$，其中 $M_i$ 取在固定的共轭类中。这些多元组要模去 $G_\IC$ 对所有 $A_i$、$B_i$、$M_i$ 的同时共轭作用。若考虑非正则穿刺点，情况就稍微复杂一些：在这样的奇点处，局部信息不止单值性；平坦截面的渐近行为还额外给出每个 $M_i$ 到斯托克斯因子与形式单值性的分解。更深入的讨论见 \cite{Witten:2007td}。

我们将在下文看到，线算子的 vev 是 $\CM$ 上的自然函数，例如上述诸矩阵之乘积的迹。这一描述使我们得以研究它们的半环关系，并与我们在第 4 节% TODO \ref{sec:Formal-line-operators}
中所作的讨论相比较。另一方面，要研究线算子与自旋特征标之间的关系，我们需要辨认出第 6 节% TODO \ref{sec:linehol}
中出现过的函数 $\CY_\gamma$。它们并不能由上述矩阵直接给出；事实上，对一般的 $G$，我们根本不知道如何构造它们（除非去解 \cite{Gaiotto:2008cd} 的积分方程）。然而对 $A_1$ 理论，我们确实知道 $\CY_\gamma$ 是什么\cite{Gaiotto:2009hg}，这一事实将在第 10 节% TODO \ref{sec:Examples}
中得到利用。

\subsubsection{同源的 Hitchin 模空间}\label{subsec:Hitchin-Isogenous}

我们在第 6.3 节% TODO \ref{sec:moreonmoduli}
中的一般讨论使我们预期：模空间不止一个 $\CM$，而是有若干个密切相关的 $\CM_\CL$，它们由线算子的族 $\CL$ 标记。这一选择在 Hitchin 系统中如何体现出来？

在定义 Hitchin 方程之解的模空间时，人们总得模去规范变换群。一种可能是使用取值为 $G$ 的单连通形式的规范变换，这给出某个模空间 $\CM_{sc}$。然而人们也可以更宽松一些，允许取值为 $G$ 的伴随形式的规范变换。这两种选择的差别是一个有限群 $\Delta$，它同构于 $C$ 上结构群为 $\CZ=Z(\tilde G)$ 的平坦丛所成的群。$\Delta$ 通过等距作用在 $\CM_{sc}$ 上，因此我们可以考虑 $\CM_{sc}$ 对 $\Delta$ 的各个子群取商。我们希望把这些商与各个 $\CM_\CL$ 对应起来。

在具体例子中，我们可以更精确地陈述 $\CL$ 的选取与 $\Delta$ 的子群之间的词典。对带正则奇异性的 $A_1$ 理论，关系特别简单：稍后我们就会看到，$\CL$ 的一个选取将对应于 $C$ 上一组闭路径 $\wp$ 的选取。我们的提法是：$\CM_\CL$ 是 $\CM_{sc}$ 对 $\Delta$ 的某个子群取商，该子群对应于沿 $\CL$ 中所有路径 $\wp$ 都具有平凡和乐的那些丛。对带非正则奇异性的 $A_1$ 理论，情况更为困难：特别地，$\CM_{sc}$ 并不是出现的最大的空间——看来还需要对所允许的规范变换施加进一步的离散限制。本文不打算作完整的分析，但会给出几个例子。

四维规范群的选取与 Hitchin 系统中允许的规范变换之间的关系，其六维动机是什么？我们没有完整的故事，但确实有一些有用的线索。尽管 Hitchin 系统背后的李代数由六维 SCFT 的选取决定，规范群的具体选取却要微妙得多。这些问题与上文第 7.1 节% TODO \ref{subsec:6d-review}
中讨论的沿圆约化时的问题类似。想来，在黎曼面上紧化时，四维规范群的选取将对应于配分函数向量空间上某种投影的选取。当我们再沿一个圆从四维约化到三维时，就没有进一步选择投影的余地了：我们对 $\CN=2$ 理论四维规范群所作的任何选取，都将决定 $C$ 上五维理论的性质。

\subsubsection{扭局部系统} \label{sec:twisted-loc}

还有第二个需要关注的整体问题。在第 6.4 节% TODO \ref{sec:twisted-vev}
中我们强调过：除了以通常方式把理论在 $S^1$ 上紧化所得的模空间 $\CM$ 之外，还存在一个先验地不同的空间 $\tCM$，它出现在以 $(-1)^{2\CJ_3}$ 作扭转的紧化中。这个空间是什么？

我们的提法是：$\tCM$ 是 $\CM$ 的一个略微扭转的版本。这里不描述 Hitchin 方程相应的扭转；我们只描述 $\tCM$ 作为复流形、在其复结构 $J^{(\zeta)}$（$\zeta\in\IC^\times$）下是什么样子。它看起来非常像 $\CM$，只不过把平坦联络换成了\emph{扭}平坦联络。$C$ 上的扭平坦 $SL(2,\IC)$-联络，是 $C$ 的单位切丛上的平坦 $SL(2,\IC)$-联络，它绕每条纤维的和乐为 $-1$。（注意，若把上一句中的 $-1$ 换成 $+1$，我们就退回到通常的平坦联络。）

如果能从 $(2,0)$ 理论的物理中直接导出这一扭转，那将是可取的。就本文的目的而言，我们满足于一个自洽性论证：如此定义的 $\tCM$ 将具有我们在第 6.4 节% TODO \ref{sec:twisted-vev}
中描述的全部性质。特别地，这一扭转很好地消解了一些棘手的符号问题。

在 $C$ 上选取一个自旋结构，给出同构 $\tCM\simeq\CM$。这是因为：固定一个自旋结构，等价于固定 $C$ 的单位切丛上一个取值为 $\IZ_2$ 的平坦局部系统，它绕每条纤维的和乐为 $-1$；与该局部系统作张量积即得所需的同构。于是这两个空间同构，但并非典范同构。一个例外是 $C$ 亏格为零的情形：此时只有一个自旋结构，因此这两个空间确实是典范同构的。

扭局部系统的模空间也曾在 \cite{MR2233852,Frenkel:2007tx} 中、出于密切相关的（但并不完全相同的）理由被考虑过。

\subsection{来自六维的线算子}

我们可以用相互作用 $(2,0)$ 理论的非阿贝尔面算子 $\IS_{\vec n}(\CR,\Sigma)$，在 $d=4,\CN=2$ 理论中定义线算子 $L_\zeta(\CR,q;\wp)$，其中 $q$ 是 $\IR^{1,3}$ 中的一条路径，$\wp\subset C$ 是一条不自交的路径。令 $\Sigma=\wp\times q$，并在第 7.2 节% TODO \ref{subsec:Twisted-C}
的扭转理论中工作。此时扭转已经破缺了 $R$ 对称性 $5=2\oplus3$，而我们希望保留 $so(3)'$ 这部分 $R$ 对称性，因此 $\vec n$ 应当是分解 $5=2\oplus3$ 中的那个二维子空间里的向量。现在取 $q=\{\vec 0\}\times\IR$，即 $\IR^{1,3}$ 中沿时间方向的一条直线。若把面算子所保持的未破缺超对称的条件——式 (7.4)% TODO \eqref{eq:surface-susy-cond}
——写出来，就会发现：只要我们允许 $\vec n$ 沿路径变化，它们实际上就是第 2 节% TODO \ref{sec:Susy-BPS-line}
中式 (2.2)% TODO \eqref{eq:Q-zeta-susy}
那种形式的四维超对称。其实这里应当说得更精确些：扭转之后，$\vec n$ 是曲线的切丛中的一个向量。所能得到的最好的「$\vec n$ 为常向量」的概念，是让 $\vec n$ 沿 $\wp$ 平移，使得 $\vec n$ 与 $\wp$ 的切向量之间的夹角 $\beta$ 保持为常数。这个选择正是保持四个超对称的正确选择，这些超对称由相位 $\zeta=e^{i\beta}$ 标记。

现在，正如我们在第 6 节% TODO \ref{sec:linehol}
中注意到的，若把时间作 Wick 转动成一个半径为 $R$ 的欧氏圆，则 vev $\langle L_\zeta(\CR,q;\wp)\rangle$（此时 $q$ 是欧氏时间圆）可以解释为超凯勒模空间 $\CM$ 上的全纯函数。利用六维的观点，我们可以给这些全纯函数一个很好的解释。

首先考虑面算子的红外极限：它由张量多重态理论中取值为 $\liet$ 的面算子写出。令
\begin{equation}\label{eq:varph-Y}
\varphi = \half\left( Y^1 + i Y^2 \right)
\end{equation}
（与 \cite{Gaiotto:2009hg} 的式 (3.31) 相同。）由于扭转，它成为 $C$ 上的一个 $(1,0)$-形式。于是对有限的 $R$，在红外极限下必须把场替换为它们沿 $S^1$ 的零模，从而
\begin{equation}
2\pi i \int_{q\times \wp} (\vec n\cdot Y) \vol(\Sigma) \longrightarrow \pi \int_\wp \left( \frac{R \varphi}{\zeta} + R \zeta \bar \varphi \right)
\end{equation}
（整体的 $i$ 因子被时间圆的欧氏延拓去掉。）类似地，$2\pi i\int_{q\times\wp}(v,B)\longrightarrow \int_\wp(v,A)$，其中 $A$ 是五维阿贝尔超 Yang--Mills 理论的一个阿贝尔规范场，$v$ 是 $\CR$ 中的一个权向量。\footnote{右端不出现 $2\pi$ 因子，是因为我们对规范场采用标准的归一化约定，使得 $F/2\pi$ 具有整周期。}

回忆：非阿贝尔 $(2,0)$ 理论在 $S^1$ 上的紧化给出五维非阿贝尔超 Yang--Mills 理论。现在 \eqref{eq:varph-Y} 定义了 $C$ 上一个取值于伴随表示的 $(1,0)$-形式 $\varphi$。鉴于上面关于阿贝尔理论的论断，自然可以猜想：非阿贝尔面算子的 vev 为
\begin{equation}\label{eq:surface-op-vev-to-5d}
\left\langle \IS_{\vec n}(\CR, q\times \wp) \right\rangle = \left\langle \Tr_{\CR} P \exp \int_\wp \left( \frac{\pi R \varphi}{\zeta} + A + \pi R \zeta \bar \varphi \right) \right\rangle,
\end{equation}
其中 $\varphi$ 与 $A$ 现在取值于规范代数为 $\lieg$ 的五维非阿贝尔 SYM 多重态，右端的 vev 是在五维 SYM 理论中取的。这无非是说：非阿贝尔面算子应当下降为五维非阿贝尔理论的 Wilson 环。左端的 vev 是 $\CM$ 上的全纯函数。此外，得益于扭转，它与 $C$ 的尺度无关，因此可以把 $C$ 取得任意大。最后，五维 SYM 理论的红外极限是自由的。这些事实强烈提示：\eqref{eq:surface-op-vev-to-5d} 中的等式是精确的，而且右端可以用其经典值代替。于是 \eqref{eq:surface-op-vev-to-5d} 的右端就是 $\CA(\zeta)$ 在 $\wp$ 上的和乐，我们最终得到本小节的主要结果：
\begin{equation}\label{eq:L-5d-Holon}
\langle L_\zeta(\CR,q;\wp) \rangle = \Tr_\CR \Hol_\wp \CA(\zeta).
\end{equation}
答案只通过 $\wp$ 的同伦类依赖于 $\wp$。这反映了如下事实：扭转已经消除了对 $C$ 上度规尺度的依赖。

若考虑如第 7.3.2 节% TODO \ref{sec:twisted-loc}
那样的、在 $S^1$ 上的扭转紧化，结果 \eqref{eq:L-5d-Holon} 只需稍作修改。对 $C$ 中任一条有向的、闭的、不自交的路径 $\wp$，都可以（用单位速度参数化 $\wp$）对应出单位切丛中一条有向闭路径 $\tilde\wp$。于是可以取扭平坦联络 $\tilde\CA$ 沿 $\tilde\wp$ 的和乐：
\begin{equation}\label{eq:L-5d-Holon-twisted}
\langle L_\zeta(\CR,q;\wp) \rangle' = \Tr_\CR \Hol_{\tilde\wp} \tilde\CA(\zeta).
\end{equation}
这给出 $\tCM$ 上的一个全纯函数。

\subsection{线算子的编目}

关于线算子的这一实现，我们再作几点评论。

\begin{enumerate}

\item 上面我们讨论的是与闭的、不自交的路径相联系的线算子。那么确有自交的路径又如何呢？表达式 $\Tr_\CR \Hol_\wp\CA(\zeta)$ 对自交路径 $\wp$ 显然有意义，但我们对六维理论中自交面算子的物理理解得还不够，无法把它们纳入讨论。人们或许还会想到 $C$ 上的「汇点」：在那里，对应于三个或更多表示 $\CR_i$ 的面算子汇合在一起。在五维 Yang--Mills 中，这类交点的数目将由 $\otimes_i\CR_i$ 中平凡表示的个数来计数：Wilson 环的端点会与相应的缠结算子缩并。$C$ 上的二维 Toda 理论（\cite{Drukker:2010jp}）中存在类似的汇点，而该理论与四维规范理论的 $S^4$ 紧化有关。这提示它们也应当是六维理论的一个特征，并给出四维理论中的线算子。我们在此不再深究。

\item 我们的构造有一些冗余。例如，\cite{Drukker:2009id} 论证道：当 ADE 群为 $A_1$ 且 $C$ 只有正则穿刺点时，沿 $C$ 上不自交路径、携带 $SU(2)$ 基础表示的面算子，已足以描述全部紫外规范理论算子：似乎并不需要更高的表示。

\item 另一方面，若考虑带非正则穿刺点的 $C$，我们迄今的讨论并没有给出全部线算子：还需要纳入涉及 $C$ 上开路径的额外算子。这里不给出这类算子的一般构造，但将在下文第 9.2 节% TODO \ref{sec:missing}
中描述它们在 $A_1$ 理论中必定呈现的样子。

\item 最后，由于互局部性的约束，这里讨论的线算子并不能同时全部引入。允许的线算子集合 $\CL$ 的选取，在这一语境下对应于对我们所允许的表示与闭路径施加某些限制。例如在 $A_1$ 理论中，我们可以自由地允许携带 $SU(2)$ 整数自旋表示的任意路径，但只能在某子集 $S$ 的路径上允许半整数自旋。局部性要求 $S$ 的选取使得 $S$ 中路径之间的相交数总是\emph{偶数}。相应的模空间 $\CM_\CL$ 是 $SU(2)$ Hitchin 系统的模空间，再模去沿 $S$ 中所有路径都具有平凡和乐的 $\IZ_2$-联络所成的群。

\end{enumerate}
